\subsubsection*{Численный пример: все шесть методов на одной функции}

\begin{idea}[title={Откуда этот материал}]
На лекции ни один из методов не был прогнан на числах. Ниже все шесть
применены к одной и той же задаче — так видно, чем они отличаются по числу
обращений к $f$ и по скорости сокращения интервала.
\end{idea}

\begin{example}\label{ex:one-dim-run}
\[
  f(x) = (x-2)^2 \ \longrightarrow\ \min , \qquad x \in [0, 5] .
\]
Функция унимодальна на $[0,5]$: она строго убывает на $[0,2]$ и строго
возрастает на $[2,5]$, так что $x_* = 2$ — единственная точка минимума.
Производная нигде не используется, только сравнения значений.

\medskip
\noindent\textbf{1) Равномерный поиск.} Возьмём $n = 5$, тогда
$h = \frac{5-0}{5} = 1$ и узлы $x_i = 0,1,2,3,4,5$:
\[
  f = 4,\ 1,\ 0,\ 1,\ 4,\ 9 .
\]
Минимум по сетке на $j = 2$ ($x_2 = 2$), новый интервал
\[
  [x_1,\, x_3] = [1,\, 3] , \qquad \text{длина } 2 ,
\]
ценой $6$ вычислений функции.

\medskip
\noindent\textbf{2) Локализация экстремума.} Тот же шаг $h = 1$. Идём слева:
\[
  i = 1: \quad f(0) = 4 > f(1) = 1 , \quad\text{но}\quad f(1) = 1 \not\le f(2) = 0 ;
\]
\[
  i = 2: \quad f(1) = 1 > f(2) = 0 \quad\text{и}\quad f(2) = 0 \le f(3) = 1 .
\]
Условия выполнены, обход прекращается: $x_* \in [1, 3]$. Результат тот же, что
в равномерном поиске, но понадобилось $4$ вычисления вместо $6$.

\medskip
\noindent\textbf{3) Дихотомия} ($\eps = 0{,}1$).

\emph{Итерация 1.} $[a_1,b_1] = [0,5]$, середина $2{,}5$:
\[
  x_1 = 2{,}4 , \quad x_2 = 2{,}6 ,
  \qquad f(x_1) = 0{,}16 , \quad f(x_2) = 0{,}36 .
\]
$f(x_1) < f(x_2)$, значит $[a_2,b_2] = [a_1,\, x_2] = [0;\ 2{,}6]$.

\emph{Итерация 2.} Середина $1{,}3$:
\[
  x_1 = 1{,}2 , \quad x_2 = 1{,}4 ,
  \qquad f(x_1) = 0{,}64 , \quad f(x_2) = 0{,}36 .
\]
Теперь $f(x_1) \ge f(x_2)$, значит $[a_3,b_3] = [x_1,\, b_2] = [1{,}2;\ 2{,}6]$.

Длины: $5 \to 2{,}6 \to 1{,}4$ за $4$ вычисления.

\medskip
\noindent\textbf{4) Бисекция.}

\emph{Итерация 1.} $[0,5]$, $L = 5$:
\[
  x_1 = 1{,}25 , \quad x_2 = 2{,}5 , \quad x_3 = 3{,}75 ,
\]
\[
  f(x_1) = 0{,}5625 , \quad f(x_2) = 0{,}25 , \quad f(x_3) = 3{,}0625 .
\]
Здесь $f(x_1) > f(x_2)$, и далее $f(x_2) \le f(x_3)$, значит берётся средняя
часть: $[a_2,b_2] = [x_1,\, x_3] = [1{,}25;\ 3{,}75]$.

\emph{Итерация 2.} $L = 2{,}5$:
\[
  x_1 = 1{,}875 , \quad x_2 = 2{,}5 , \quad x_3 = 3{,}125 ,
\]
\[
  f(x_1) = 0{,}015625 , \quad f(x_2) = 0{,}25 , \quad f(x_3) = 1{,}265625 .
\]
Теперь $f(x_1) < f(x_2)$, значит $[a_3,b_3] = [a_2,\, x_2] = [1{,}25;\ 2{,}5]$.

Длины: $5 \to 2{,}5 \to 1{,}25$ — ровно вдвое за итерацию, как и обещано.

\medskip
\noindent\textbf{5) Золотое сечение}, $\alpha = \frac{3-\sqrt5}{2} \approx 0{,}381966$.

\emph{Итерация 1.} $[a_1,b_1] = [0,5]$:
\[
  \lambda_1 = 0 + \alpha\cdot 5 \approx 1{,}90983 ,
  \qquad
  \mu_1 = 5 - \alpha\cdot 5 \approx 3{,}09017 ,
\]
\[
  f(\lambda_1) \approx 0{,}00813 , \qquad f(\mu_1) \approx 1{,}18847 .
\]
$f(\lambda_1) < f(\mu_1)$, значит $[a_2,b_2] = [a_1,\, \mu_1] = [0;\ 3{,}09017]$,
причём $\mu_2 = \lambda_1 \approx 1{,}90983$ — значение уже известно.

\emph{Итерация 2.} Вычисляется \emph{одна} новая точка:
\[
  \lambda_2 = a_2 + \alpha(b_2 - a_2) \approx 1{,}18034 ,
  \qquad f(\lambda_2) \approx 0{,}67079 .
\]
Теперь $f(\lambda_2) \ge f(\mu_2)$, значит
$[a_3,b_3] = [\lambda_2,\, b_2] = [1{,}18034;\ 3{,}09017]$.

Длины: $5 \to 3{,}09017 \to 1{,}90983$ за $3$ вычисления функции — по одному на
итерацию после первой.

\medskip
\noindent\textbf{6) Метод Фибоначчи}, $n = 5$:
\[
  F_0 = 1,\ F_1 = 1,\ F_2 = 2,\ F_3 = 3,\ F_4 = 5,\ F_5 = 8 .
\]

\emph{Итерация 1} ($k=1$): коэффициенты
$\frac{F_{n-k-1}}{F_{n-k+1}} = \frac{F_3}{F_5} = \frac38$ и
$\frac{F_{n-k}}{F_{n-k+1}} = \frac{F_4}{F_5} = \frac58$, поэтому
\[
  \lambda_1 = 0 + \tfrac38\cdot 5 = 1{,}875 ,
  \qquad
  \mu_1 = 0 + \tfrac58\cdot 5 = 3{,}125 ,
\]
\[
  f(\lambda_1) = 0{,}015625 , \qquad f(\mu_1) = 1{,}265625 .
\]
$f(\lambda_1) < f(\mu_1)$, значит $[a_2,b_2] = [0;\ 3{,}125]$ и
$\mu_2 = \lambda_1 = 1{,}875$.

\emph{Итерация 2} ($k=2$): коэффициенты $\frac{F_2}{F_4} = \frac25$ и
$\frac{F_3}{F_4} = \frac35$:
\[
  \lambda_2 = \tfrac25\cdot 3{,}125 = 1{,}25 ,
  \qquad
  \mu_2 = \tfrac35\cdot 3{,}125 = 1{,}875 \quad\text{(совпало, как и должно)} ,
\]
\[
  f(\lambda_2) = 0{,}5625 , \qquad f(\mu_2) = 0{,}015625 .
\]
$f(\lambda_2) \ge f(\mu_2)$, значит $[a_3,b_3] = [1{,}25;\ 3{,}125]$.

Длины: $5 \to 3{,}125 \to 1{,}875$.
\end{example}

\begin{idea}
Совпадение $\mu_2 = \lambda_1 = 1{,}875$ в методе Фибоначчи — это и есть
переиспользование точки, доказанное выше; поэтому на второй итерации
вычисляется только $f(\lambda_2)$.

Сравнение по длине интервала на одно вычисление функции (все методы стартовали
с $[0,5]$):
\begin{itemize}[nosep,left=0pt]
  \item равномерный поиск: $2$ за $6$ вычислений;
  \item локализация: $2$ за $4$;
  \item дихотомия: $1{,}4$ за $4$;
  \item бисекция: $1{,}25$ за $4$ (с переиспользованием середины — за $2$ новых);
  \item золотое сечение: $1{,}90983$ за $3$;
  \item Фибоначчи: $1{,}875$ за $3$.
\end{itemize}
Фибоначчи обходит золотое сечение именно потому, что число шагов задано
заранее и коэффициенты подобраны под него; при $n \to \infty$ разница исчезает.
\end{idea}
